\pg{337} 61. Corollarios. --- Dos principios expostos resulta:

I. O ente \emph{necessario}, assim como na sua essencia \emph{ideal} contem a razão da sua existencia \emph{ideal}, assim tambem na sua essencia \emph{real} contem a razão da sua existencia \emph{real}; porque a existencia pertence a mesma ordem, á qual pertence a \emph{essencia} (\textsuperscript{1}).

II. Acertada a existencia \emph{real} de um só ente \emph{contingente}, deduz-se legitimamente a existencia do ente \emph{necessario}; pois o ente \emph{contingente} não tem em si a razão da sua existencia \footnote{(1) O attributo segue sempre a condição do sujeito e por isso encontra-se na mesma ordem, em que se encontra o sejeito. --- Encontra-se o sujeito na ordem \emph{ideal}? O attributo deve encontrar-se na ordem \emph{ideal}. --- Encontra-se o sujeito na ordem \emph{real}? O attributo tambem deve encontrar-se na ordem \emph{real}. Ao ente \emph{necessario} convem \emph{essencialmente} a \emph{existencia}. Portanto, onde se encontra o ente \emph{necessario}; ahi deve encontrar-se a existencia. Mas se o ente \emph{necessario} existisse \emph{exclusivamente} na intelligencia, não poderia ter senão uma existencia \emph{ideal}. Por isso, quando se quer deduzir da \emph{idea} do ente \emph{necessario} a sua existencia \emph{real}, commette-se um sophisma. --- A existencia pertence \emph{essencialmente} á \emph{idea} de ente \emph{necessario}. Por isso, se esse ente existe, \emph{essencial} e \emph{necessariamente} existe. Mas é necessario provar que effectivamente o ente \emph{necessario} existe. Assim, pertencendo a razão á essencia do homem, segue-se que, se existe o homem, existe uma creatura, a qual é essencial e necessariamente racional; mas do facto de ser o homem essencial e necessariamente racional não segue que essa creatura existe na realidade. Essa existencia deve provar-se com argumentos fundados, não nas ideas, mas nos factos.}.\footnote{(2) A demonstração da existencia do ente \emph{necessario} deve fundar-se nos factos. Existem no mundo entes \emph{contingentes}? A experiencia interna e externa diz que existem; porque ha no mundo entes que começam e acabam, que estão sujeitos a mudanças, e por isso apresentam um caracter de \emph{contingencia}. Se existem os entes \emph{contingentes}, deve existir o ente \emph{necessario}. O \pgn{338}proprio \emph{Kant} escreve: « Desde que existe alguma coisa, não podemos deixar de concluir que existe tambem um ente \emph{necessario} ». --- Nem se diga que não podemos concluir da existencia do ente \emph{contingente} a existencia do ente \emph{necessario}, por não haver nenhuma relação entre o \emph{contingente} e o \emph{necessario}. Entre o \emph{contingente} e o \emph{necessario} existe, como vimos, a relação, que existe entre o effeito e a causa; e por isso, assim como da existencia do effeito deduzimos legitimamente a existencia da causa respectiva, assim tambem da existencia do ente \emph{contingente} deduzimos a existencia do ente \emph{necessario}.}
